% Generated by report/tools/tables.py from the result files of the repository. Do not edit.
\begin{table}[tbp]
\centering
\footnotesize
\setlength{\tabcolsep}{3.2pt}
\caption[Plus petite différence détectable]{Plus petite différence que le sondage pouvait détecter : pour chaque critère, partant de la part de 2018 et des effectifs réels, la plus petite hausse et la plus petite baisse, en points, que le test exact de Fisher bilatéral détecte avec 80~\% de puissance, au seuil de 5~\% puis de 5~\%/11 (Bonferroni sur les onze critères v1) ; « aucune » : aucune variation possible n'atteint cette puissance. À droite, la puissance qu'avait le test pour l'écart effectivement vu. \textbf{En gras}, les quatre critères dont la différence passe la correction de Holm (tableau~\ref{tab:sondage-pineau}) : tous dépassent la différence détectable. Ailleurs, l'absence de changement détecté n'est pas une preuve d'absence : la hausse de Q7 (+12 points) n'avait que 18~\% de chances d'être vue au seuil de 5~\%. Source : \file{study/survey\_extra.json} (\file{study/survey\_extra.py}).}
\label{tab:difference-detectable}
\begin{tabularx}{\linewidth}{@{}l>{\raggedright\arraybackslash}Xrrrrrrrr@{}}
\toprule
 & & & & \multicolumn{2}{c}{Au seuil de 5~\%} & \multicolumn{2}{c}{\makecell{Bonferroni\\{}(5~\%/11)}} & \multicolumn{2}{c}{\makecell{Puissance de\\{}l'écart vu}} \\
\cmidrule(lr){5-6}\cmidrule(lr){7-8}\cmidrule(lr){9-10}
Code & Critère & 2018 (\%) & écart vu & hausse & baisse & hausse & baisse & 5~\% & Bonf. \\
\midrule
Q2 & Réseaux de neurones employés & 79 & +19 & +18 & −28 & aucune & −36 & 88 & 56 \\
Q3 & Tous les hyperparamètres de la méthode & 55 & +17 & +27 & −29 & +34 & −37 & 38 & 11 \\
Q4 & Tous les hyperparamètres des références & 35 & −7 & +29 & −24 & +37 & −30 & 9 & 1 \\
\textbf{Q5} & \textbf{Code en lien} & \textbf{31} & \textbf{+43} & \textbf{+29} & \textbf{−23} & \textbf{+38} & \textbf{−28} & \textbf{99} & \textbf{93} \\
Q6 & Choix des hyperparamètres décrit & 26 & +22 & +30 & −21 & +38 & −24 & 57 & 24 \\
Q7 & Évaluation sur conditions tenues à part & 33 & +12 & +29 & −24 & +37 & −29 & 18 & 4 \\
Q8 & Test statistique appliqué & 6 & +6 & +22 & aucune & +30 & aucune & 9 & 1 \\
\midrule
\textbf{X2} & \textbf{Incertitude rapportée} & \textbf{63} & \textbf{+26} & \textbf{+25} & \textbf{−29} & \textbf{+31} & \textbf{−38} & \textbf{83} & \textbf{52} \\
X3 & Seuls les meilleurs essais gardés & 19 & −5 & +27 & −17 & +36 & aucune & 8 & 1 \\
X4 & Références réglées avec le même soin & 13 & +4 & +26 & aucune & +35 & aucune & 5 & 1 \\
\textbf{X7} & \textbf{Calcul chiffré} & \textbf{27} & \textbf{+31} & \textbf{+28} & \textbf{−21} & \textbf{+38} & \textbf{−25} & \textbf{88} & \textbf{59} \\
\midrule
Y1 & Mesure de performance définie & 79 & −1 & +18 & −28 & aucune & −37 & 4 & 0 \\
Y2 & Barres d'erreur définies & 83 & +3 & aucune & −33 & aucune & −43 & 4 & 0 \\
Y3 & Plusieurs tâches (au moins 2) & 90 & −8 & aucune & −24 & aucune & −33 & 13 & 2 \\
Y4 & Une tâche perdue montrée & 79 & +14 & +20 & −31 & aucune & −40 & 36 & 9 \\
Y5 & Origine du code des références & 46 & −5 & +29 & −27 & +36 & −34 & 7 & 1 \\
Y6 & Exclusions dites & 29 & +3 & +28 & −22 & +39 & −26 & 5 & 0 \\
Y7 & Version des bancs d'essai & 23 & +3 & +28 & −19 & +37 & aucune & 5 & 0 \\
Y8 & Test de généralisation & 23 & +11 & +28 & −19 & +37 & aucune & 20 & 4 \\
\textbf{Y9} & \textbf{Sensibilité à un hyperparamètre} & \textbf{44} & \textbf{+36} & \textbf{+30} & \textbf{−27} & \textbf{+37} & \textbf{−34} & \textbf{95} & \textbf{79} \\
Y10 & Versions des bibliothèques & 2 & +6 & +19 & aucune & +26 & aucune & 9 & 1 \\
Y11 & Type d'incertitude nommé & 83 & +10 & aucune & −33 & aucune & −43 & 21 & 4 \\
Y12 & Outil statistique postérieur à 2018 & 2 & +12 & +19 & aucune & +26 & aucune & 44 & 13 \\
\bottomrule
\end{tabularx}
\end{table}
